\begin{center}\begin{tikzpicture}[font=\small,>=stealth,line join=round]
\draw[gray!40,rounded corners=3pt] (0,0) rectangle (14.2,8.3);\node at (7.1,7.700000000000001) {例2：两条奇函数相乘，得到左右同高};
\begin{scope}[shift={(4.1,2.161111111111111)},x=1.6666666666666665cm,y=1.0222222222222221cm]\draw[->,gray] (-1.8,0)--(1.8,0) node[right]{$x$};\draw[->,gray] (0,-0.5)--(0,4) node[above]{$y$};\draw[blue,thick,solid,domain=-1.5:1.5,samples=180] plot(\x,{(\x)^3*sin(deg(\x))});\end{scope}\node[anchor=west,align=left,text width=5.0cm] at (8.4,4.0) {$F=x^3\sin x$\\[10pt]两因子都变号\\[10pt]乘积不变号\\[10pt]图像关于 $y$ 轴对称};
\node[text width=12.6cm,align=center] at (7.1,.65) {这里讨论的是函数的对称性；不能套用整数“奇数乘奇数”的说法。};
\end{tikzpicture}\end{center}
